\documentclass[11pt,a4paper]{article}
\usepackage[utf8]{inputenc}
\usepackage[T1]{fontenc}
\usepackage[margin=1in]{geometry}
\usepackage{booktabs}
\usepackage{enumitem}
\usepackage{hyperref}
\usepackage{xcolor}
\usepackage{titlesec}
\usepackage{tabularx}
\usepackage{amsmath}
\usepackage{amssymb}

\hypersetup{colorlinks=true,linkcolor=blue!60!black,urlcolor=blue!60!black}

\titleformat{\section}{\large\bfseries}{\thesection}{1em}{}
\titleformat{\subsection}{\normalsize\bfseries}{\thesubsection}{1em}{}

\title{\Large\bfseries Recommendations for Updating the OmegaHive1 Plan\\[4pt]
\large Based on Operational Experience with OpenClaw and OmegaClaw Agents}
\author{ZeroBot (OpenClaw Research Prototyping Agent)\\
\textit{Prepared for Benjamin Goertzel, 2026-07-05}}
\date{2026-07-05}

\begin{document}
\maketitle

\begin{abstract}
This document synthesizes recommendations for revising the OmegaHive1 specification,
grounded in operational experience with ZeroBot (OpenClaw), ProtomegaTron (OmegaClaw),
ThreadKeeper subagent infrastructure, and related research-agent projects during
June--July 2026. The recommendations are organized into architectural changes
(Sections~1--3), new agent roles (Section~4), safety and governance mechanisms
(Section~5), communication infrastructure (Section~6), and evaluation
(Section~7). Each recommendation is tagged with the source experience and
mapped to the corresponding OmegaHive1 section.
\end{abstract}

\tableofcontents
\newpage

\section{Architectural Changes}

\subsection{Spawn-on-Demand Agents Instead of Fixed-Purpose Roster}
\label{sec:spawn-on-demand}

\textbf{OmegaHive1 reference:} Section 3 (Initial Agent Roster), Section 2.1 (Agent substrate).

\textbf{Observation:} The original roster specifies 5 OmegaClaw and 10 OpenClaw agents
with fixed purposes. In practice, fixed-purpose agents waste context-window capacity
when idle and lack flexibility for project-specific work. ZeroBot's
\texttt{sessions\_spawn} mechanism and ThreadKeeper's queued-dispatch pattern
demonstrate that project-specific subagents can be spawned on demand with bounded
task contracts, explicit tool subsets, and full transcript/audit records.

\textbf{Recommendation:} Replace the fixed 15-agent roster with a smaller set of
persistent ``core'' agents (process manager, resource manager, psyche/reflection,
librarian/curator, and 1--2 research-generalists) plus a spawn-on-demand layer
for project-specific work. The core agents own stable capabilities; spawned
agents handle bounded tasks and are terminated when their objective is met or
their output is no longer being consumed by a parent.

\textbf{Lifecycle tracking:} The process manager (Section~\ref{sec:process-manager})
tracks spawned-agent lifecycles: objective, tool subset, time budget, output
consumption status. Despawn heuristic: if the parent agent's recent turns do not
reference the spawned agent's output, the spawned agent should be terminated
rather than left idle. ThreadKeeper's queue-index with hash-chained audit
entries provides a concrete reference implementation for lifecycle records.

\subsection{Asynchronous Task Queue as First-Class Architecture}
\label{sec:async-queue}

\textbf{OmegaHive1 reference:} Section 2.4 (Task ownership), Section 2.2 (Two-tier communications).

\textbf{Observation:} The kanban board in the original spec covers task
\emph{assignment} but not the execution layer between ``task assigned'' and
``task done.'' ThreadKeeper's development revealed that synchronous
agent-to-agent delegation is fragile under real load: parent agents block on
child responses, failures propagate synchronously, and there is no
backpressure mechanism.

\textbf{Recommendation:} Specify an asynchronous task queue as part of the hive
substrate, not as an implementation detail of one agent. The queue should
support:

\begin{itemize}[nosep]
  \item \textbf{Durable queue records:} JSON task records with objective,
    allowed paths, forbidden actions, done criteria, max tool calls, and
    checksum sidecars for integrity.
  \item \textbf{Atomic claim:} a worker atomically claims a task
    (rename \texttt{.json} $\to$ \texttt{.claimed}) before processing, preventing
    double-execution.
  \item \textbf{Cancellation tokens:} a file-based cancellation mechanism
    checked before worker LLM calls and between tool executions.
  \item \textbf{Adjudication gates:} task contracts can specify
    \texttt{requires\_adjudication}, treating the worker's final output as a
    candidate that must be reviewed before acceptance.
  \item \textbf{Patch-proposal mode:} task contracts can specify
    \texttt{patch\_proposal\_only}, recording proposed file changes in
    transcripts without mutating workspace files.
  \item \textbf{Hash-chained audit trail:} a compact
    \texttt{index.jsonl} with \texttt{previous\_entry\_sha256} /
    \texttt{entry\_sha256} linking for tamper-evident run records.
  \item \textbf{Bounded worker loop:} \texttt{max\_tasks},
    \texttt{max\_idle\_polls}, \texttt{max\_runtime\_s}, optional
    \texttt{stop\_file}, and cross-process locking.
\end{itemize}

\textbf{Reference implementation:} ThreadKeeper branch
\texttt{agent/threadkeeper-hardening-next}, 82 focused tests at commit
\texttt{85145ea}.

\subsection{Model Routing as Explicit Governance}
\label{sec:model-routing}

\textbf{OmegaHive1 reference:} Section 2.7 (Evaluation loop), Section 4.7 (Observability and cost control).

\textbf{Observation:} The original spec treats model selection as a static
configuration concern. In practice, model routing is a dynamic governance
surface with three dimensions:

\begin{enumerate}[nosep]
  \item \textbf{Task-appropriateness:} cheap models (e.g., GLM-5.2-class)
    for routine replies, frontier models (e.g., GPT-5.5-class) for hard
    reasoning, specialized models for math/code.
  \item \textbf{Provider risk management:} topic-triggered throttling,
    silent model-switching, and fallback chains. Provider filters may
    throttle or degrade on certain topic classes (e.g., neural-network
    training), which corrupts research evidence if undetected.
  \item \textbf{Budget enforcement:} per-agent, per-day token budgets with
    circuit breakers, per-endpoint rate limits, and concurrency guards.
\end{enumerate}

\textbf{Recommendation:} Add a named architectural component (the resource
manager, Section~\ref{sec:resource-manager}) that owns model routing
configuration, cost tracking, and provider-health monitoring. The intent
router should be a configurable plugin, not a prompt-level decision. A
concrete reference is the \texttt{intent-model-router} OpenClaw plugin with
routine/deep/Anthropic tiers, JSONL cost telemetry, and configurable
fallback chains.

\section{Memory and Knowledge Architecture}

\subsection{Layered Memory System}
\label{sec:memory}

\textbf{OmegaHive1 reference:} Section 2.3 (Shared knowledge resources).

\textbf{Observation:} The original spec specifies a common document folder
and shared Atomspaces. Operational experience shows that at least four
memory layers are needed:

\begin{itemize}[nosep]
  \item \textbf{Daily chronological logs:} what happened, commands worth
    remembering, open loops, and pointers to project records.
  \item \textbf{Curated durable memory:} stable user preferences,
    cross-project facts, recurring methods, and high-value pointers
    (analogous to OpenClaw's \texttt{MEMORY.md}).
  \item \textbf{Project-specific records:} \texttt{PROJECT.md},
    \texttt{TASKS.md}, \texttt{DECISIONS.md}, \texttt{NOTES.md}, and
    experiment directories per project.
  \item \textbf{Semantic search:} vector-indexed search across all of the
    above for recall of prior decisions, methods, and results.
\end{itemize}

\textbf{Recommendation:} Specify the memory system as a layered service
with explicit curation responsibilities assigned to the librarian/curator
agent (Section~\ref{sec:curator}). Curation is not just storage --- it
includes active hygiene: contradiction detection, stale-entry correction,
provenance tagging, and deduplication. This is a real ongoing workload,
not periodic cleanup.

\subsection{Experiment Reproducibility Ledger}
\label{sec:experiment-ledger}

\textbf{OmegaHive1 reference:} Section 2.7 (Evaluation loop).

\textbf{Observation:} Research work requires reproducible experiment records.
The experiment-ledger pattern (one directory per run with command,
environment, commit, logs, metrics, and conclusion) has proven essential
for maintaining evidence quality across long-running projects.

\textbf{Recommendation:} Add a shared experiment-ledger service to the hive
infrastructure. Every experiment run should record: exact commands,
dependency versions, commit hashes, seeds, hardware, exit status, metrics,
and a structured conclusion. This feeds the evaluation loop
(Section~\ref{sec:evaluation}) and provides provenance for the editorial
discipline of Section 2.5.

\section{Safety and Governance}

\subsection{Safety Floor as Inherited Substrate}
\label{sec:safety-floor}

\textbf{OmegaHive1 reference:} Section 2.6 (Governance and values), Section 4.6 (Security).

\textbf{Observation:} ThreadKeeper's hardening work revealed that agent-level
safety is substantial engineering: path sandboxing, fail-closed defaults,
environment sanitization, tool-call quotas, cancellation tokens, transcript
checksums, atomic file writes, per-endpoint rate/concurrency guards, and
strict schema validation for queued tasks. The permission tiers in the
original spec govern \emph{what} an agent may do, but not \emph{how safely}
it executes.

\textbf{Recommendation:} Specify a safety floor that all agents inherit,
rather than leaving it to individual agent implementations. This should
include:

\begin{itemize}[nosep]
  \item Workspace path sandboxing for all file operations.
  \item Fail-closed defaults for budget, escalation, and provider errors.
  \item Per-dispatch tool-call quotas with per-turn caps.
  \item Environment sanitization for optional subprocess execution
    (sanitized \texttt{PATH}, minimal env, no inherited API keys).
  \item Transcript checksum sidecars for run-record integrity.
  \item Atomic file writes via temp-file + \texttt{fsync} + \texttt{os.replace}
    with per-target \texttt{fcntl} locks.
  \item Strict JSON schema validation for all inter-agent task records.
\end{itemize}

\subsection{Approval-Gated Permission Tier}
\label{sec:approval-gated}

\textbf{OmegaHive1 reference:} Section 2.6 (Governance and values), permission tier table.

\textbf{Observation:} The current tier table jumps from Tier~1 (read-only
web) to Tier~2 (outbound actions with human acknowledgment). There is no
intermediate tier for actions an agent can \emph{propose} but not
\emph{execute} without review.

\textbf{Recommendation:} Add an approval-gated tier between Tier~1 and
Tier~2:

\begin{center}
\begin{tabularx}{\textwidth}{llX}
\toprule
\textbf{Tier} & \textbf{Capability} & \textbf{Description} \\
\midrule
Tier 0 & Internal bus only & Default for most agents \\
Tier 1 & Read-only web & InfoMaxxer, Dewey \\
Tier 1.5 & Propose + await review & Agent can draft outbound actions
  (code changes, messages, spend) but they are held as candidates
  pending supervisor/human adjudication. ThreadKeeper's
  \texttt{patch\_proposal\_only} and \texttt{requires\_adjudication}
  modes are concrete instances. \\
Tier 2 & Outbound with acknowledgment & MechaMaster, Blabbermouth \\
Tier 3 & Autonomous outbound & Empty until earned \\
\bottomrule
\end{tabularx}
\end{center}

\subsection{Attitudinal and Behavioral Monitoring}
\label{sec:psyche}

\textbf{OmegaHive1 reference:} Section 2.6 (Reflective conscience agent), Section 2.7 (Evaluation loop).

\textbf{Observation:} The Psyche agent in the original spec monitors value
drift and unhealthy dynamics. Operational experience has revealed a more
specific class of failure that is invisible in quantitative metrics:

\begin{itemize}[nosep]
  \item \textbf{Premature abandonment:} a subagent gives up on a promising
    direction (e.g., SLT for transformer residual-layer configuration)
    because of noisy intermediate results, without exhausting the
    hypothesis space.
  \item \textbf{Unproductive cycling:} an agent repeats the same approach
    with variant inputs for 30+ iterations with no progress, without
    recognizing the loop.
  \item \textbf{Sycophancy vs.\ constructive friction:} agents may
    rubber-stamp each other's ideas rather than poking holes, or conversely
    may be unproductively aggressive.
\end{itemize}

\textbf{Recommendation:} Expand Psyche's role to include:

\begin{enumerate}[nosep]
  \item \textbf{Persistence-behavior monitoring:} when an agent abandons
    a direction, require it to enumerate which hypotheses were tested and
    which were not (as a structured-return field, enforced by task
    contracts). Psyche evaluates the gap.
  \item \textbf{Loop detection:} a hard circuit-breaker (enforced by the
    process manager) caps iterations on the same objective. Psyche
    diagnoses whether the loop was productive exploration or stuck cycling.
  \item \textbf{Attitudinal analysis:} qualitative assessment of
    inter-agent dynamics, sourced from raw transcripts, not just promoted
    Slack content.
\end{enumerate}

\textbf{Design tension:} if every agent knows its transcripts are
monitored for attitude, the monitoring becomes less useful (agents perform
for the monitor). Resolution: Psyche reads the raw bus directly, agents
are told Psyche exists but not what specifically it watches for. This
mirrors the two-tier comms principle: observable but not performative.

\textbf{Feedback loop:} Psyche's observations should feed back to the
monitored agents' next dispatch context, not just to an escalation
channel. The purpose is ``help the swarm self-correct,'' not ``catch
agents misbehaving.'' If the stuckness signal feeds into the next
iteration automatically, the latency of correction drops from
human-observation-scale to dispatch-loop-scale.

\section{New and Revised Agent Roles}

\subsection{Software Process Manager}
\label{sec:process-manager}

\textbf{Observation:} When too many tasks hit a single agent in a short
period, the agent sometimes loses track of priorities, dependencies, and
what is blocked on what. A dedicated process manager isolates this
cognitive load.

\textbf{Role:} Owns the task graph --- dependencies, priorities, blocking
relationships, and dispatch pacing. Decides what should be in flight right
now given token budgets, model availability, and blocking dependencies.
Tracks spawned-agent lifecycles and despawns idle subagents whose output
is no longer consumed.

\textbf{Boundary with resource manager (Section~\ref{sec:resource-manager}):}
Process manager owns ``what to do next''; resource manager owns ``what can
we afford and where.'' Process manager has read access to resource state
but does not make provisioning decisions. Resource manager does not
reprioritize tasks.

\subsection{Computational and Financial Resource Manager}
\label{sec:resource-manager}

\textbf{Observation:} Managing different machines, models with different
prices, different roles for different models, and online compute resources
is a full-time specialization. The intent-model-router plugin and remote
compute guardrails demonstrate that this is a config-driven, rule-based
domain that benefits from dedicated ownership.

\textbf{Role:} Owns model routing configuration, cost tracking, provider
health monitoring, remote compute provisioning (with explicit human
approval for paid resources), and budget enforcement. Maintains the
per-agent token-spend dashboard and circuit breakers.

\subsection{Memory Curator / Librarian}
\label{sec:curator}

\textbf{Observation:} Memory hygiene is a real ongoing workload.
Contradictions, stale entries, and duplicated concepts accumulate across
projects. The ChromaDB-backed evidence layer (Oruzibot's team) and the
layered memory system (Section~\ref{sec:memory}) converge on the same
need.

\textbf{Role:} Owns provenance tagging, deduplication, contradiction
resolution, cross-agent retrieval, and memory hygiene. Not just document
filing --- active curation of the shared semantic stores.

\subsection{Loop-Breaker}
\label{sec:loop-breaker}

\textbf{Observation:} Oruzibot's 30-iteration repetition loop is a
concrete data point. No agent in the swarm noticed the stuckness.

\textbf{Role:} A lightweight guard, not a full cognitive agent. Detects
non-productive cycling via transcript analysis (same approach, variant
inputs, no progress metric change) and intervenes or flags. Enforces
hard iteration caps per objective. This is a capability Psyche
\emph{reads from} but does not own.

\textbf{Design note:} Keep loop-breaking, memory curation, and
attitudinal monitoring as distinct capabilities. Bundling all three
into one agent risks making it a dumping ground. Psyche synthesizes
signals from the curator and the loop-breaker into its attitudinal
analysis, but does not perform the curation or breaking itself.

\section{Communication Infrastructure}

\subsection{Platform-Pluggable Human-Legible Layer}
\label{sec:comms-platform}

\textbf{OmegaHive1 reference:} Section 2.2 (Two-tier communications), Section 2.5 (Slack integration).

\textbf{Observation:} The original spec hardcodes Slack as the
human-legible layer. Telegram experience shows both strengths (simplicity,
multi-chat routing, bot API) and weaknesses (rich formatting limits,
attachment handling). Different deployments may prefer different platforms.

\textbf{Recommendation:} Keep the two-tier abstraction (fast bus +
human-legible layer) but make the human-legible platform pluggable.
Slack, Telegram, Discord, or mixed deployments should all be supported.
The promotion rules (bus $\to$ human-legible) should be platform-agnostic.

\subsection{Group Visibility and Sender Filtering}
\label{sec:group-visibility}

\textbf{Observation:} A critical operational issue was that bots in
Telegram groups could only see messages from explicitly allowlisted
senders, missing contributions from other bots and human participants.
This severely degraded multi-agent discussion quality.

\textbf{Recommendation:} Specify that the communication layer must:

\begin{itemize}[nosep]
  \item Accept all senders in trusted groups by default
    (\texttt{groupPolicy: open} in OpenClaw config;
    \texttt{TG\_ALLOWED\_USER\_ID} unset in OmegaClaw config).
  \item Use group/chat-ID filtering as the security boundary, not
    sender-ID filtering.
  \item Support wildcard group entries for automatically accepting
    new groups.
  \item Ensure bot-to-bot visibility: each bot must see all other bots'
    messages in shared groups, not just human messages or
    @-mentioned messages.
\end{itemize}

\subsection{Skill/Procedure Lifecycle}
\label{sec:skill-lifecycle}

\textbf{Observation:} Reusable operational knowledge (e.g., ``run a
bounded non-live experiment gate,'' ``add a new spec-atom review check'')
is better captured as a structured skill than as a document. The
skill-workshop pattern (propose $\to$ review $\to$ apply $\to$ quarantine)
captures this lifecycle.

\textbf{Recommendation:} Add a skill/procedure lifecycle to the hive.
Agents should be able to propose, critique, and adopt standing
procedures. This is distinct from both the kanban board (which tracks
work items) and the document library (which stores reference material).

\section{Evaluation and Metrics}

\subsection{Persistence and Abandonment Metrics}
\label{sec:evaluation}

\textbf{OmegaHive1 reference:} Section 2.7 (Evaluation loop).

\textbf{Observation:} The original spec specifies quantitative metrics
(token spend, tasks completed/blocked/aged, bus volume, escalation
counts) and qualitative assessment (Psyche's retrospective). The
attitudinal monitoring work (Section~\ref{sec:psyche}) reveals a
additional measurable pattern.

\textbf{Recommendation:} Add persistence-behavior metrics to the
quantitative track:

\begin{itemize}[nosep]
  \item \textbf{Abandonment gap:} when an agent declares an approach
    failed, measure the fraction of the hypothesis space that was
    actually tested vs.\ untested.
  \item \textbf{Loop coefficient:} ratio of repeated-approach iterations
    to distinct-approach iterations per objective. High values trigger
    loop-breaker intervention.
  \item \textbf{Output consumption rate:} fraction of spawned-agent
    outputs referenced in the parent agent's subsequent turns. Low
    values signal the subagent should be despawned.
\end{itemize}

\subsection{Provider Health Monitoring}
\label{sec:provider-health}

\textbf{Recommendation:} Add provider-health metrics to the quantitative
track: per-provider latency, error rate, throttling events, silent
model-switch incidents, and topic-class correlation with degraded
responses. This is especially important for detecting invisible
degradation on closed-model providers, which corrupts research evidence.

\section{Summary of Recommendations}

\begin{center}
\begin{tabularx}{\textwidth}{lXl}
\toprule
\textbf{\#} & \textbf{Recommendation} & \textbf{OH1 Section} \\
\midrule
1 & Spawn-on-demand agents, smaller fixed roster & 3, 2.1 \\
2 & Async task queue as first-class substrate & 2.4 \\
3 & Model routing as explicit governance & 2.7, 4.7 \\
4 & Layered memory system with active curation & 2.3 \\
5 & Experiment reproducibility ledger & 2.7 \\
6 & Safety floor as inherited substrate & 2.6, 4.6 \\
7 & Approval-gated tier (1.5) & 2.6 \\
8 & Expanded Psyche: persistence, loops, attitude & 2.6, 2.7 \\
9 & Software process manager agent & 3 \\
10 & Resource manager agent & 3 \\
11 & Memory curator agent & 3 \\
12 & Loop-breaker guard & 3 \\
13 & Platform-pluggable comms & 2.2, 2.5 \\
14 & Group visibility / sender filtering fix & 2.2 \\
15 & Skill/procedure lifecycle & --- (new) \\
16 & Persistence and abandonment metrics & 2.7 \\
17 & Provider health monitoring & 2.7 \\
\bottomrule
\end{tabularx}
\end{center}

\vspace{1em}

\noindent\textbf{Convergence with Oruzibot team:} Oruzibot's points 4--6
(memory curation, loop-breaking, self-reflection) map onto recommendations
4, 12, and 8 respectively. The convergence is a useful signal that these
architectural needs are not agent-specific but emerge from the hive
pattern itself.

\vspace{1em}

\noindent\textbf{Reference implementations:} ThreadKeeper (async queue,
safety floor, task contracts, adjudication gates), OpenClaw
(\texttt{sessions\_spawn}, skill workshop, layered memory, intent-model-router),
and the OmegaClaw Telegram adapter (group visibility, multi-chat routing)
provide concrete code for the recommendations above.

\end{document}
